% !TeX program = xelatex
% CS336 Spring 2026 / Lecture 02 / Personal review notes
% Self-contained source. Compile twice with XeLaTeX for cross-references.
\documentclass[UTF8,fontset=fandol,zihao=-4,a4paper]{ctexart}
\usepackage[margin=22mm,headheight=16pt,headsep=7mm,footskip=11mm]{geometry}
\usepackage{amsmath,amssymb,booktabs,tabularx,array,xcolor,enumitem,fancyhdr}
\usepackage[T1]{fontenc}
\usepackage{textcomp,listings}
\usepackage[unicode,colorlinks=true,linkcolor=blue!55!black,urlcolor=blue!55!black]{hyperref}
\definecolor{accent}{HTML}{205F99}
\definecolor{muted}{HTML}{536170}
\definecolor{soft}{HTML}{F2F6FA}
\setlength{\parindent}{2em}
\setlength{\parskip}{0.35em}
\linespread{1.13}
\setlist[itemize]{leftmargin=1.7em,itemsep=0.25em,topsep=0.35em}
\setlist[enumerate]{leftmargin=2em,itemsep=0.35em,topsep=0.35em}
\setlength{\emergencystretch}{2em}
\renewcommand{\arraystretch}{1.3}
\newcolumntype{Y}{>{\raggedright\arraybackslash}X}
\newcolumntype{P}[1]{>{\raggedright\arraybackslash}p{#1}}
\ctexset{
 section={format=\Large\bfseries\color{accent},beforeskip=0.7em,afterskip=0.45em},
 subsection={format=\normalsize\bfseries,beforeskip=0.65em,afterskip=0.25em}
}
\pagestyle{fancy}
\fancyhf{}
\fancyhead[L]{\small\color{muted}CS336 · 第 02 课}
\fancyhead[R]{\small\color{muted}PyTorch 与资源核算}
\fancyfoot[L]{\footnotesize\color{muted}Spring 2026 · 个人复习资料}
\fancyfoot[R]{\small\thepage}
\renewcommand{\headrulewidth}{0.35pt}
\lstset{basicstyle=\ttfamily\small,backgroundcolor=\color{soft},
 frame=single,rulecolor=\color{accent!20},columns=fullflexible,
 keepspaces=true,breaklines=true,showstringspaces=false,
 xleftmargin=0.6em,xrightmargin=0.6em,aboveskip=0.65em,belowskip=0.65em}
\newcommand{\focus}[1]{\begin{quote}\small\color{accent}\bfseries #1\end{quote}}
\newcommand{\tagline}[1]{\noindent{\small\color{muted}#1}\par}
\newcommand{\tok}[1]{\texttt{#1}}
\newcommand{\spacechar}{\textvisiblespace}
\newcommand{\nextpage}{\clearpage}
\hypersetup{pdftitle={CS336 第02课复习：PyTorch 与资源核算},pdfauthor={个人学习笔记},linktoc=all}
\begin{document}

\begin{center}
{\small\color{muted}STANFORD CS336 · LANGUAGE MODELING FROM SCRATCH}\\[0.55em]
{\LARGE\bfseries 第 02 课复习资料}\\[0.3em]
{\large PyTorch 与资源核算}
\end{center}

\setcounter{tocdepth}{1}
\tableofcontents
\clearpage

\tagline{阅读主线：文本变成张量之后，怎样判断模型能否装下、要算多久，以及究竟慢在哪里？}
\section{先定位：从“怎样表示文本”走向“怎样核算资源”}\label{sec:overview}

\subsection{接上第一课：表示方案最终会变成资源开销}
第一课把文本变成 token ID，随后由嵌入表变成向量。第二课继续追问：这些向量、模型参数及训练时的中间结果要占多少显存？运算需要多少计算？为什么 GPU 峰值算力很高，实际训练仍然可能很慢？

\focus{本课的核心问题：给定模型、数据和硬件，先算清“装不装得下”，再估计“算多久”，最后找出“算力还是搬运限制了速度”。}
PyTorch 是表达计算的工具，资源核算是贯穿全课的思维方式。只会写出正确结果，还不足以判断一个方案能否扩展到大模型。

\subsection{三本账：容量、工作量与速度}
\par\noindent\begin{tabularx}{\linewidth}{@{}P{23mm}P{34mm}Y@{}}
\toprule
问题 & 基本量与单位 & 核算方法\\
\midrule
装不装得下？ & 存储量，byte & 元素个数 $\times$ 每元素字节数；再汇总同时存活的张量。\\
总计算量？ & 运算量，FLOPs & 从矩阵形状计数；把前向、反向及额外重算区分开。\\
耗时下界？ & 吞吐率、带宽 & 计算量除以 FLOP/s，搬运量除以 byte/s，判断哪一项更慢。\\
\bottomrule
\end{tabularx}
显存\textbf{容量}大，表示能装更多；显存\textbf{带宽}高，表示每秒能搬更多。两者不是同一指标。总 FLOPs 少，也不自动意味着实际耗时短。

\subsection{连续例子：同一个小网络，逐步把账补完整}
沿用官方讲义的等宽网络：输入 $X\in\mathbb R^{B\times d}$，每层权重 $W_\ell\in\mathbb R^{d\times d}$，依次执行
\[
H_0=X,\qquad H_\ell=\operatorname{ReLU}(H_{\ell-1}W_\ell),
\qquad \ell=1,\ldots,L.
\]
无偏置、各层不共享权重，参数量为 $N=Ld^2$。正文固定用 $B=4,d=8,L=3$：共有 192 个参数，每层输出 32 个元素。后面依次为它补上形状、字节、FLOPs、梯度与优化器状态。

\tagline{复习路线：第 \ref{sec:tensors}--\ref{sec:training} 节建立理解；第 \ref{sec:review} 节一页速查；第 \ref{sec:quiz} 节独立自测。正文依据正式与录像版讲义，来源和差异说明见第 \ref{sec:sources} 节。}

\nextpage
\section{接上主线：张量里到底存了什么，占多少空间？}\label{sec:tensors}
\tagline{本节目标：看到 shape、dtype、device，就能开始核算；分清精度与范围。}

\subsection{先认清三个属性：形状、数据类型、位置}
张量是多维数组；本课的 rank 表示\textbf{轴的数量}，不是线性代数中的矩阵秩。形状为 $(4,8,2)$ 的张量有 3 个轴、64 个元素。数据、参数、梯度、激活和优化器状态，都可用张量存储。
\begin{lstlisting}
x = torch.zeros(4, 8)
x.shape          # (4, 8)
x.numel()        # 32
x.dtype          # torch.float32, under the usual default
x.element_size() # 4 bytes
x.device         # cpu, under the usual default
\end{lstlisting}
对普通稠密张量，逻辑数据量为
\[
M_{\text{tensor}}=\left(\prod_i s_i\right)q,
\qquad s_i=\text{第 }i\text{ 个轴的长度},\quad q=\text{每元素字节数}.
\]
例如 $(4,8)$ 的 fp32 数据为 $4\times8\times4=128$ byte；改为 bf16 则为 64 byte。多个张量若共享底层存储，不能把各自逻辑大小重复算作独立分配；实际显存还有分配器与临时缓冲区开销。

\subsection{从小例子走到 Transformer 的形状}
讲义展示 $(B,S,H,d_h)=(32,16,16,64)$：批大小、序列长度、头数、每头维度。共有 $524\,288$ 个元素，fp32 占 2 MiB，bf16 占 1 MiB。

这里 $d_h$ 是每头维度；后文 $d$ 表示整个隐藏维度，勿混用。统一单位也很重要：$1\,\mathrm{GB}=10^9$ byte，$1\,\mathrm{GiB}=2^{30}$ byte。

\subsection{为什么同样 16 bit，fp16 和 bf16 不一样？}
浮点数可理解为符号、指数和有效数字三部分。\textbf{指数}主要决定能覆盖的数量级；\textbf{有效数字位数}主要决定同一数量级内分得多细。
\par\noindent\begin{tabularx}{\linewidth}{@{}P{17mm}P{12mm}P{25mm}Y@{}}
\toprule
类型 & byte & 指数 / 尾数 & 需要记住的取舍\\
\midrule
fp32 & 4 & 8 / 23 & 范围大、分辨率较高，存储成本也高。\\
fp16 & 2 & 5 / 10 & 比 fp32 省空间，但范围较窄，易溢出或下溢。\\
bf16 & 2 & 8 / 7 & 指数范围接近 fp32，但相邻可表示数更疏。\\
\bottomrule
\end{tabularx}
表内为位数，三种类型另有 1 个符号位。讲义用 $10^{-8}$ 演示：转为 fp16 会成为 0，bf16 仍非零；这\textbf{不意味着} bf16 具有 fp32 的精度，也不保证训练必然稳定。

\nextpage
\subsection{混合精度：不同用途的张量，不必用同一种精度}
大矩阵乘法希望省显存并利用低精度硬件；长期累积的优化器统计量则需要更稳健的表示。讲义用“参数、激活、梯度取 bf16，优化器状态取 fp32”做资源估算。\textbf{这是一个明确的核算配置，不是所有训练程序的默认布局。}

原始混合精度训练论文还讨论 fp32 主权重副本和 loss scaling；若实现保留主权重，必须额外计入显存。loss scaling 先放大损失以减少小梯度下溢，再在更新前还原梯度；它不能修复所有数值问题。

\textbf{补充澄清：}PyTorch 的 autocast 按运算选择计算类型；它不会把上下文中的一切张量永久转成 bf16，也不能据此断言参数、梯度和优化器状态都是低精度。
\begin{lstlisting}
# Assuming CUDA is available, x and w are fp32 CUDA tensors.
with torch.amp.autocast("cuda", dtype=torch.bfloat16):
    y = x @ w   # eligible matmul may use bf16
    z = torch.zeros(4, 8, device="cuda")
    # z still uses the default dtype; autocast is not a global cast
\end{lstlisting}
做真实预算时，直接检查各类张量的 dtype，再乘元素数；不要从“开启 AMP”反推出固定的每参数字节数。

\subsection{把张量放到 GPU：位置也是计算约定的一部分}
CPU 内存和 GPU 显存是不同位置。普通运算的输入与权重需在兼容设备上；搬运数据也需要时间。
\begin{lstlisting}
if torch.cuda.is_available():
    y = x.to("cuda:0")
    z = torch.zeros(32, 32, device="cuda:0")
\end{lstlisting}
\tok{x.to(...)} 返回目标张量，需保存返回值。录像版还查询 GPU 数量、设备属性以及 \tok{memory\_allocated()}：两个独立的 $(32,32)$ fp32 GPU 张量，逻辑数据合计 $8192$ byte。

\subsection{再往下压位宽：fp8 与 fp4 的启示}
讲义进一步介绍 FP8 的 E4M3、E5M2：前者分辨率较好、后者范围较大；按讲义采用的格式，最大有限值分别为 448 与 57344。它们需要结合缩放等数值策略，而非直接替换所有 fp32 数值。

NVFP4 把单个数值编码压到 4 bit，并引入分块缩放。\textbf{“4 bit”只描述数值编码本身}，总存储还包括缩放因子等信息；也依赖对应硬件与库支持。本课掌握动机和取舍即可，不要求实现低精度训练配方。

\focus{第一本账的原则：先列出实际存在的张量，再按它们各自的元素数与 dtype 计费；“模型用了低精度”不是完整的显存预算。}

\nextpage
\section{理解操作：用 einops 把形状与计算说清楚}\label{sec:einops}
\tagline{本节目标：能读懂 einsum、reduce、rearrange，并把形状接到后面的成本计数。}

\subsection{einsum：保留哪些轴，沿哪些轴求和？}
下文的 \tok{einsum}、\tok{reduce}、\tok{rearrange} 均从 \tok{einops} 导入。
对 $X\in\mathbb R^{B\times d}$、$W\in\mathbb R^{d\times k}$，线性变换为
\[
Y_{bo}=\sum_{i=1}^{d}X_{bi}W_{io},\qquad Y\in\mathbb R^{B\times k}.
\]
\begin{lstlisting}
y = einsum(x, w, "batch in, in out -> batch out")
\end{lstlisting}
出现在输出中的 \tok{batch,out} 被保留；未出现在输出中的 \tok{in} 被求和。轴名是形状说明，不会自动识别数据的真实语义；相同名称的轴仍须有正确对应关系。

讲义的成对相似度例子：$X:(B,S_1,d)$，$Y:(B,S_2,d)$。
\begin{lstlisting}
z = einsum(x, y,
    "batch seq1 hidden, batch seq2 hidden -> batch seq1 seq2")
\end{lstlisting}
结果形状为 $(B,S_1,S_2)$，对每对位置沿 hidden 轴做内积。即使 $S_1=S_2$，也应使用不同轴名，避免误写成同位置逐元素相乘。省略号 \tok{...} 可表达共同的前导维度。

\subsection{reduce：明确写出汇总方式}
\begin{lstlisting}
x = torch.ones(2, 3, 4)
y = reduce(x, "batch seq hidden -> batch seq", "sum")
\end{lstlisting}
reduce 可以求和、均值、最大值等；输出中消失的轴被归约。\tok{sum} 与 \tok{mean} 的缩放不同，会影响损失与梯度的大小，不能只检查输出形状。

\subsection{rearrange：拆轴、换轴、合轴，不等于做乘加}
沿用讲义 $(3,8)$ 的输入，把总宽度 $8=2\times4$ 拆成两个头，每头作用同一个 $(4,4)$ 权重，再合回去。
\begin{lstlisting}
x = rearrange(x, "seq (heads din) -> seq heads din", heads=2)
x = einsum(x, w, "seq heads din, din dout -> seq heads dout")
x = rearrange(x, "seq heads dout -> seq (heads dout)")
\end{lstlisting}
这里两个头\textbf{共享同一个} $W$。若每头有独立参数，权重应多一个 heads 轴，参数量和写法都会改变。rearrange 保持元素总数；有时能返回视图，有时需要复制，不能一概当作零显存、零耗时。

\focus{读一段张量代码时，先写输入形状，再标出求和轴，最后写输出形状。形状一旦明确，显存与 FLOPs 都有了起点。}

\nextpage
\section{核心例子：从一次矩阵乘法推到整次训练}\label{sec:flops}
\tagline{本节目标：推导 $2Bdk$、反向的两次矩阵乘法，以及 $6NT$ 近似的适用范围。}

\subsection{FLOPs 是工作量，FLOP/s 才是速度}
一次浮点加法或乘法通常各计 1 FLOP；融合乘加虽可由一条硬件指令执行，惯例仍计 2 FLOPs。FLOPs 表示总工作量，FLOP/s 表示每秒工作量，TFLOP/s 中的 T 表示 $10^{12}$。

对 $Y=XW$，$X:(B,d)$、$W:(d,k)$、$Y:(B,k)$，每个输出需 $d$ 次乘法与 $d-1$ 次加法，所以
\[
F_{\text{exact}}=Bk(2d-1)\approx 2Bdk.
\]
两个同形矩阵相加则只需 $Bk$ 次加法。大尺寸时矩阵乘法通常主导 FLOPs；这句话尚未说明谁主导\textbf{耗时}，后面还需计入搬运。

回到小网络：$B=4,d=8,L=3,N=192$，每层前向约 $2\times4\times8^2=512$ FLOPs，三层约 1536 FLOPs。若严格只数乘加，则为 $3\times4\times8\times15=1440$；ReLU 与损失另计。\textbf{精确计数与量级估算不能混成同一个数字。}

\subsection{先用一个标量损失理解 backward 在求什么}
官方例子为 $x=(1,2,3)$、$w=(1,1,1)$，预测 $p=x^\top w=6$，目标为 5：
\[
\mathcal L=\tfrac12(p-5)^2=\tfrac12,
\qquad \nabla_w\mathcal L=(p-5)x=(1,2,3).
\]
\begin{lstlisting}
w = torch.tensor([1., 1., 1.], requires_grad=True)
loss = 0.5 * (x @ w - 5).pow(2)
loss.backward()   # w.grad == tensor([1., 2., 3.])
\end{lstlisting}
\tok{requires\_grad=True} 告诉自动微分系统跟踪相关运算。默认主要在需要梯度的叶子张量上保留 \tok{.grad}；中间量通常不会保留可直接读取的梯度，可用 \tok{retain\_grad()} 调试。\tok{.grad} 为 None 不等于这一步没有参与反向计算。

\subsection{为什么一个线性层的反向约为前向的两倍？}
设上游梯度 $G=\partial\mathcal L/\partial Y\in\mathbb R^{B\times k}$，则
\[
\frac{\partial\mathcal L}{\partial X}=GW^\top\quad(B\times d),
\qquad
\frac{\partial\mathcal L}{\partial W}=X^\top G\quad(d\times k).
\]
两次矩阵乘法各约 $2Bdk$ FLOPs，合计约 $4Bdk$。反向既要把梯度传给前层，又要求当前参数的梯度；后一式还解释了为什么反向常需要保留前向输入 $X$。

\nextpage
\subsection{从每层相加，得到 $6N$ 这条估算规则}
等宽网络有 $N=Ld^2$ 个参数。忽略激活函数、损失、优化器等较小项，并按各层都求输入与参数梯度估算：
\[
F_{\text{forward}}\approx2BN,\qquad
F_{\text{backward}}\approx4BN,\qquad
F_{\text{train step}}\approx6BN.
\]
小网络得到 $6\times4\times192=4608$ FLOPs。这是统一口径的\textbf{粗估}；若首层输入不需要梯度，实际会省去相应矩阵乘法。检查点重算还会额外增加工作量。

语言模型中，每条序列有多个 token。将 $T$ 定义为\textbf{整次训练累计处理的 token 数}，则常用
\[
\boxed{F_{\text{training}}\approx6NT.}
\]
这里 $T$ 不是单条序列长度，$N$ 也应与所计入的主要稠密参数运算匹配。该近似适合矩阵投影等项占主导的情形；长上下文注意力、词表投影、稀疏激活的 MoE、检查点重算等，需要另外核算，不能无条件套用全部参数量。

\subsection{回到开场问题：70B 模型训练 15T token 要多久？}
沿用讲义 H100 SXM 的估算口径：稠密 bf16 Tensor Core 峰值
$P=1979/2=989.5\,\mathrm{TFLOP/s}$，1024 张卡，假设有效利用率 $u=0.5$。
\[
F=6\times70\times10^9\times15\times10^{12}
=6.3\times10^{24}\ \mathrm{FLOPs},
\]
\[
t\approx\frac{F}{1024\times989.5\times10^{12}\times0.5}
\approx1.244\times10^7\ \mathrm{s}\approx143.9\ \text{天}.
\]
这不是工期承诺：假设利用率已能概括主要运行损耗，没有另算停机、数据处理、评估等时间；换硬件或计算类型时，要换对应峰值。稀疏峰值不能直接用于普通稠密矩阵乘法。

\subsection{MFU：离所选硬件峰值还有多远？}
若时间窗口内模型的有效计算量为 $F_{\text{model}}$，测得耗时为 $t$，则
\[
\mathrm{MFU}=\frac{F_{\text{model}}/t}{n_{\text{GPU}}P_{\text{peak}}}.
\]
分子分母必须使用一致的精度、稠密／稀疏和统计范围。重算是实际执行的额外工作，一般不当作模型有效 FLOPs 加进 MFU 分子。小网络也不能期待接近大矩阵的峰值。

\textbf{计时补充：}CUDA 执行通常异步。应预热、重复测量，并在计时边界同步，或使用 CUDA events；只计 Python 提交运算的时间会低估耗时。算完 FLOPs 后，还要解释实际吞吐为什么达不到峰值，这就是下一节。

\nextpage
\section{理解速度：为什么算得少，也可能被显存带宽卡住？}\label{sec:intensity}
\tagline{本节目标：把 FLOPs 与搬运字节数放在一起，用 roofline 判断瓶颈。}

\subsection{一次运算有两项成本：计算与搬运}
即使数据已经在 GPU 上，仍需从显存 HBM 读取，送到计算单元，最后写回结果。本节讨论的是\textbf{GPU 内部显存与计算单元之间的搬运}，不是把 CPU 到 GPU 的拷贝混进来。

设运算量为 $F$、HBM 搬运量为 $Q$ byte、计算峰值为 $P$ FLOP/s、显存带宽为 $\beta$ byte/s，则
\[
t_{\text{compute}}=\frac{F}{P},\qquad
t_{\text{memory}}=\frac{Q}{\beta},\qquad
t\gtrsim\max\!\left(\frac{F}{P},\frac{Q}{\beta}\right).
\]
右边是理想条件下的\textbf{耗时下界}。讲义假设计算与搬运充分重叠，用两者较大项近似时间；真实运行还受启动、同步、缓存、调度等因素影响。

\subsection{算术强度：每搬 1 byte，能做多少计算？}
定义工作负载的算术强度 $I$，以及硬件的算力／带宽比 $I_*$：
\[
I=\frac{F}{Q},\qquad I_*=\frac{P}{\beta},\qquad
\text{单位均为 FLOPs/byte}.
\]
若 $I<I_*$，搬运相对更慢，称为 memory-bound（带宽受限）；若 $I>I_*$，计算相对更慢，称为 compute-bound（计算受限）。注意，这里的 memory-bound \textbf{不是显存容量不足}。

沿讲义的 H100 数值 $P=989.5\times10^{12}$、$\beta=3.35\times10^{12}$：
\[
I_*\approx295.4\ \mathrm{FLOPs/byte}.
\]
这个阈值对应选定的稠密 bf16 计算峰值；逐元素运算未必在相同计算单元上执行，其细致分析要用相应运算的峰值。这里按课堂简化模型建立量级直觉。

\subsection{roofline：斜线是带宽天花板，横线是计算天花板}
由时间下界可推出有效吞吐上界：
\[
P_{\text{achievable}}\leq\min(P,\beta I),\qquad
\frac{P_{\text{achievable}}}{P}\leq\min\!\left(1,\frac{I}{I_*}\right).
\]
横轴是算术强度、纵轴是吞吐：左侧随 $I$ 增大而上升，右侧达到计算天花板，拐点为 $I_*$。真实 MFU 还会更低，不能把这个\textbf{理想上界}直接当作实测值。

\focus{带宽受限时，优先减少搬运或提高数据复用；计算受限时，再考虑减少计算或提高矩阵乘法效率。优化前，先判断处在哪一侧。}

\nextpage
\subsection{五个例子：为什么大矩阵乘法特别适合 GPU？}
沿用讲义的 bf16 假设：每个输入／输出元素 2 byte，输入从 HBM 读一次、输出写一次，暂不计额外中间缓冲。ReLU 比较按一次基本操作计；GELU 的 20 次操作只是课堂粗估。
\par\noindent\begin{tabularx}{\linewidth}{@{}P{34mm}P{29mm}P{31mm}Y@{}}
\toprule
运算 & 工作量 $F$ & 搬运量 $Q$ & 强度 $I$\\
\midrule
ReLU，$n$ 个数 & $n$ & $4n$ & $1/4$\\
GELU，$n$ 个数 & 约 $20n$ & $4n$ & 约 $5$\\
长度 $n$ 的内积 & $2n-1$ & $4n+2$ & 趋近 $1/2$\\
$n\times n$ 矩阵乘向量 & $n(2n-1)$ & $2n^2+4n$ & 趋近 $1$\\
两个 $n\times n$ 矩阵相乘 & $n^2(2n-1)$ & $6n^2$ & 约 $n/3$\\
\bottomrule
\end{tabularx}
前三类到矩阵向量乘法，强度不会随规模显著增长；大矩阵乘法则能反复复用每个输入元素。$n=1024$ 时，矩阵乘法强度约 $341.2>295.4$，在此理想模型中进入计算受限区。

对于 $n=2^{20}$ 的独立 ReLU，搬运 $4\,194\,304$ byte，理想搬运时间约 $1.25\,\mu\mathrm{s}$，按课堂峰值估的计算时间约 $1.06\,\mathrm{ns}$。两者相差很大；实际计时还会有启动等成本。GELU 虽多做运算，也可能主要等同样的数据搬运，因此减少 FLOPs 未必带来等比例加速，\textbf{但不能据此保证二者实测一样快}。

\subsection{用同一个公式连接训练、批处理和生成}
对 bf16 的 $(B,d)\times(d,k)$ 矩阵乘法：
\[
I\approx\frac{2Bdk}{2(Bd+dk+Bk)}
=\frac{Bdk}{Bd+dk+Bk}.
\]
当 $B$ 相对 $d,k$ 很小时，读权重 $dk$ 占主导，$I\approx B$。单条序列逐 token 解码常接近矩阵向量乘法，读一遍权重只服务很少 token；把更多 token 合批，能摊薄权重搬运，增加复用。

训练中的投影通常同时处理多个 token，矩阵更大，更容易利用算力。这里不能概括为“所有训练都计算受限，所有推理都带宽受限”：预填充、大批量解码、注意力和 KV cache 都会改变账本。

\subsection{这也解释了为什么融合运算有用}
\textbf{补充解释：}若矩阵乘法后紧接激活，分开执行可能把中间结果先写回 HBM，再读回来。融合可减少这次往返。数学结果和主要 FLOPs 近似不变，但 $Q$ 下降，速度仍可能提升。更深入的 kernel 与 GPU 优化留待后续课。

\nextpage
\section{把账算完整：训练状态、优化器与两种省显存方法}\label{sec:training}
\tagline{本节目标：不漏算梯度和优化器状态；理解梯度累积与激活检查点分别减少什么。}

\subsection{从一个训练循环理解张量的生命周期}
\begin{lstlisting}
optimizer.zero_grad(set_to_none=True)
for x, target in batches:
    pred = model(x)
    loss = criterion(pred, target)  # a scalar loss
    loss.backward()
    optimizer.step()
    optimizer.zero_grad(set_to_none=True)
\end{lstlisting}
前向保留反向需要的信息；backward 累加参数梯度；step 用梯度和状态更新参数；zero\_grad 清理梯度。\textbf{step 通常不自动清梯度}。set\_to\_none 设梯度为 None，分配器仍可能保留缓存，驱动看到的占用未必立即下降。

官方网络用 \tok{nn.Parameter} 注册权重，用 \tok{nn.ModuleList} 注册多层模块，使参数遍历、设备移动和状态保存能找到它们。普通输入与中间激活不会仅因是张量就自动成为可训练参数。

\subsection{优化器为什么还要保存一份“历史”？}
AdaGrad 对每个参数维护累计梯度平方。按讲义实现的稳定项位置，逐元素更新为
\[
v_t=v_{t-1}+g_t^2,\qquad
\theta_{t+1}=\theta_t-\eta\frac{g_t}{\sqrt{v_t+\varepsilon}}.
\]
经常累积较大梯度的坐标，后续更新会被缩小。$v_t$ 与参数形状相同，因此即使参数不变多，也多出一份状态。AdaGrad 是\textbf{累加}平方，不要与 RMSProp 的指数滑动平均混淆。

动量法保留平滑梯度，RMSProp 保留梯度平方的指数平均，Adam 同时保存一阶与二阶统计；AdamW 增加解耦权重衰减。本课重在状态成本，不展开完整配方。

\subsection{小网络的显存账：按用途逐行相加}
假设参数和梯度为 bf16，优化器状态为 fp32；用 $N=192$、$B=4,d=8,L=3$。按课堂的简化激活模型估算：
\par\noindent\begin{tabularx}{\linewidth}{@{}P{34mm}Y P{21mm}@{}}
\toprule
项目 & 计算式 & 字节数\\
\midrule
参数 & $2N$ & 384\\
梯度 & $2N$ & 384\\
AdaGrad 状态 & $4N$ & 768\\
激活（粗估） & $2BLd$ & 192\\
\midrule
合计（AdaGrad） & $8N+2BLd$ & 1728\\
合计（Adam） & $12N+2BLd$ & 2496\\
\bottomrule
\end{tabularx}

\nextpage
\subsection{从小网络推广：每参数 12 byte 是带前提的}
上一页 Adam 的 $12N$ 来自：2 byte 参数、2 byte 梯度、两份各 4 byte 的统计量。它\textbf{尚未包含}激活、fp32 主权重副本、输入数据、临时空间与分配器开销。激活的 $2BLd$ 也只是量级模型，不是 PyTorch 自动微分保存所有张量的精确峰值。

开场的 8 张 80 GB H100，在理想完整切分模型状态的条件下，仅按 $12N$ 算：
\[
N_{\max}\lesssim\frac{8\times80\times10^9}{12}
\approx53.3\times10^9\quad\text{个参数}.
\]
这是\textbf{只算模型状态的理想容量上界}。若普通数据并行在每卡复制完整状态，就不能相加使用 8 张卡显存；单卡相同口径上界约为 6.67B。若另有 fp32 主权重，状态变成 $16N$，8 卡聚合上界便降到 40B，激活等还需继续扣除。

\subsection{梯度累积：把一个大批次拆成多个小批次}
设有效批大小为 $B$，拆为 $K$ 个相等的 micro-batch，每个大小 $b=B/K$。逐个做前向、反向，累积梯度，全部完成后只更新一次。这样可及时释放前一个 micro-batch 的计算图，激活峰值从约 $2BLd$ 变为 $2bLd$；参数、梯度和优化器状态并未因此缩小。
\begin{lstlisting}
# K equal-sized micro-batches; criterion returns their mean loss.
optimizer.zero_grad(set_to_none=True)
for x, target in micro_batches:
    loss = criterion(model(x), target) / K
    loss.backward()
optimizer.step()
optimizer.zero_grad(set_to_none=True)
\end{lstlisting}
若每个 micro-batch 损失已取均值，需要除以 $K$ 才与大批次均值对应；大小不等时按样本数加权，语言模型按参与损失的有效 token 数加权。不能在每个 micro-batch 后清梯度或更新参数，否则变成另一种训练过程。

讲义例子 $B=64,d=1024,L=16$：简化 bf16 激活量为 2 MiB；拆成 4 个大小 16 的 micro-batch 后约为 0.5 MiB。总样本工作量近似不变，但小矩阵的算术强度和硬件效率可能下降。

\focus{梯度累积解决“同一次更新想看更多数据，但激活放不下”的问题。它通常不减少模型状态，也不承诺提升速度。}
\tagline{等价性前提：损失可按样本／token 分解且参数在累积期间固定；批统计、随机性及浮点求和次序可能带来差异。梯度裁剪等操作通常应在累积完成后执行。}

\nextpage
\subsection{激活检查点：少存一些，需要时重新算出来}
反向需要前向中间量，例如 $\partial\mathcal L/\partial W=X^\top G$ 中的 $X$。常规方式把这些信息保存到反向；activation checkpointing（也叫 gradient checkpointing 或 rematerialization）则只保留部分边界，反向时从最近边界重做前向，恢复缺失激活。

它保存的是\textbf{反向所需的激活检查点}，不是训练中保存到磁盘、用于恢复任务的模型 checkpoint。两者解决的问题不同。

对各层激活大小相同的 $L$ 层链，设每隔 $k$ 层保存一个边界。反向处理一段时，既需要所有边界，也需要暂存当前段的中间量，因此简化峰值为
\[
M_{\text{activation}}=O\!\left(Bd\left(\frac{L}{k}+k\right)\right),
\qquad k\approx\sqrt L\ \Longrightarrow\ O(Bd\sqrt L).
\]
字节数还需乘 dtype 宽度。以 $L=16$ 为例，每 4 层一个边界，避免同时保留全部层内部激活；这给出的是量级关系，不是精确节省比例。

\par\noindent\begin{tabularx}{\linewidth}{@{}P{39mm}P{33mm}Y@{}}
\toprule
策略（简化链模型） & 激活空间量级 & 额外重算工作量\\
\midrule
保存所有层所需激活 & $O(L)$ & 无需重算前向。\\
逐层从头恢复、只留常数状态 & $O(1)$ & 总重算可达 $O(L^2)$。\\
每约 $\sqrt L$ 层保存边界 & $O(\sqrt L)$ & 总额外重算为 $O(L)$。\\
\bottomrule
\end{tabularx}
表中省略了 $Bd$ 与 dtype 等共同因子。最后一行的 $O(L)$ 指\textbf{所有段加起来的重算}；不是把段长 $O(\sqrt L)$ 当作全网重算量。实际选择取决于模块大小、可用显存和重算开销。

\textbf{PyTorch 接口补充：}可用 \tok{torch.utils.checkpoint.checkpoint} 包裹需要重算的函数；按当前文档应明确选择 \tok{use\_reentrant=False}。被包裹函数的重算需与原前向一致；随机状态、副作用或可变全局状态要遵守接口约定。课堂“每层包一次”的例子和“每隔 $\sqrt L$ 层分段”的理论策略也不是同一种分段方式。

\subsection{两种方法怎样共同服务主线？}
梯度累积缩小\textbf{同时处理的数据量}；激活检查点减少\textbf{同时保留的中间结果}。前者通常维持近似总 FLOPs，后者明确增加计算。它们可组合使用，但都不消除模型状态成本。

\focus{优化闭环：先列出显存最大项，再选方法；改完重新计数 FLOPs 与搬运，并用实际峰值和计时验证。能放下、算得少、跑得快，是三个相关但不同的目标。}

\nextpage
\section{一页速查：把本课压回一条主线}\label{sec:review}
\tagline{快速复习时先读这一页；无法解释某条结论时，回到对应章节看推导。}

\subsection{用六句话重建本课}
\begin{enumerate}
\item 文本变成张量之后，要同时核算显存容量、总运算量与数据搬运，才能判断训练方案是否可行。
\item 张量逻辑数据量等于元素数乘 dtype 字节数；参数、梯度、优化器状态与激活应分别列账。
\item einops 显式写出轴的含义：einsum 做乘积与求和，reduce 做归约，rearrange 做拆合和重排。
\item 线性层前向约 $2Bdk$，反向两次矩阵乘法约 $4Bdk$，从而得到常用训练粗估 $6NT$。
\item 算术强度 $I=F/Q$ 决定算力和带宽谁更紧；大矩阵乘法通过复用输入提高强度。
\item 梯度累积减少 micro-batch 激活峰值，检查点用额外重算换激活空间；两者都有适用条件与代价。
\end{enumerate}

\subsection{五条公式与一个连续例子}
\par\noindent\begin{tabularx}{\linewidth}{@{}P{58mm}Y@{}}
\toprule
记忆点 & 对应含义与边界\\
\midrule
$M=(\prod_i s_i)q$ & 稠密张量逻辑字节数；共享存储与分配开销另计。\\
$F_{\text{matmul}}\approx2Bdk$ & $(B,d)\times(d,k)$；前向乘加估算。\\
$F_{\text{train}}\approx6NT$ & $N$ 为对应参数量，$T$ 为累计 token 数；不含额外重算等。\\
$I=F/Q,\ I_*=P/\beta$ & 同为 FLOPs/byte；$I<I_*$ 时理想模型带宽受限。\\
$t\gtrsim\max(F/P,Q/\beta)$ & 时间下界；吞吐上界为 $\min(P,\beta I)$。\\
$B=4,d=8,L=3$ & $N=192$；前向约 1536 FLOPs，训练约 4608；简化 AdaGrad 显存 1728 byte。\\
\bottomrule
\end{tabularx}

\subsection{易错点：看到这些说法，要能立即纠正}
\begin{itemize}
\item \textbf{“bf16 和 fp32 一样精确。”} 范围接近，不等于有效数字位数相同。
\item \textbf{“有 8 卡，就能直接相加显存。”} 取决于状态是否切分；复制式数据并行不成立。
\item \textbf{“低 FLOPs 就一定快。”} 还要看字节搬运、形状和启动等开销。
\item \textbf{“memory-bound 就是显存不够。”} 此处是带宽受限，容量另算。
\item \textbf{“累积梯度只要多次 backward。”} 还需正确缩放损失、控制更新与清梯度时机。
\item \textbf{“检查点免费省显存。”} 激活减少，但多了前向重算；MFU 核算也需统一口径。
\end{itemize}

\nextpage
\section{自测：合上前文，检查自己是否真的理解}\label{sec:quiz}
\tagline{建议先独立作答，再看第 \ref{sec:answers} 节。每题都尝试说出计算前提与单位。}

\subsection{概念题：检查因果关系}
\begin{enumerate}[label=\textbf{Q\arabic*.}]
\item \textbf{三本账。}同一个模型换上容量更大、但带宽和算力相同的 GPU，为什么可能能训练更大 batch，却不保证单次运算更快？
\item \textbf{精度与存储。}形状 $(32,16,16,64)$ 的张量有多少元素？fp32、bf16 分别占多少 MiB？为什么 bf16 能表示 $10^{-8}$，却不能据此说它与 fp32 一样精确？
\item \textbf{形状判断。}输入 $X:(2,3,4)$、$Y:(2,5,4)$，执行 \tok{einsum} 模式 \tok{batch seq1 hidden, batch seq2 hidden -> batch seq1 seq2}，输出形状与求和轴是什么？与沿 hidden 求均值有何不同？
\end{enumerate}

\subsection{手算题：检查操作是否准确}
\begin{enumerate}[label=\textbf{Q\arabic*.},start=4]
\item \textbf{前向和反向。}$X:(4,8)$、$W:(8,8)$。用 $2Bdk$ 口径估计前向 FLOPs；写出两项反向矩阵公式及各自形状。如果该层输入不需要梯度，会少做哪一项？
\item \textbf{完整显存账。}三层无偏置、不共享权重的等宽网络，$B=4,d=8,L=3$。参数和梯度用 bf16，AdaGrad 状态用 fp32，激活按 $2BLd$ byte 粗估。参数量与总显存是多少？改成两份 fp32 状态的 Adam 后呢？
\item \textbf{算术强度。}bf16 的 $1024\times1024$ 方阵乘法，假设两个输入各读一次、输出写一次。估算 $F,Q,I$，与 $I_*=295.4$ 比较。若换成矩阵向量乘法，为什么强度会大幅下降？
\end{enumerate}

\subsection{迁移题：检查是否会判断取舍}
\begin{enumerate}[label=\textbf{Q\arabic*.},start=7]
\item \textbf{利用率与工期。}用 $6NT$ 估计 70B 模型、15T token 的总计算量。1024 张卡、每卡峰值 $989.5$ TFLOP/s、MFU 为 0.5，要多少天？若 MFU 变为 0.25，其余不变，时间如何变化？
\item \textbf{梯度累积。}有效 batch 为 64，分成 4 个等大的 micro-batch，各自损失都已取均值。应怎样缩放损失，何时 step 和 zero\_grad？哪类显存减少，哪类不减少？
\item \textbf{激活检查点。}等宽 $L$ 层链每隔 $k$ 层保留边界，为什么激活峰值含 $L/k+k$ 两项？怎样选 $k$？为什么“总重算为 $O(\sqrt L)$”不正确？它能否解决参数本身就装不下的问题？
\end{enumerate}

\focus{掌握标准：能从形状算字节和 FLOPs，解释反向的两次矩阵乘法，用算术强度判断瓶颈，并根据显存最大项选择优化方法。}
\tagline{延伸观察：任选一段现有训练代码，列出每类张量的形状、dtype 与生命周期，再对照实测峰值。先解释差值来自哪里，不急于追求一个看似精确的总数。}

\nextpage
\section{参考答案与资料来源}\label{sec:answers}
\tagline{答案用于核对关键点，不要求逐字背诵；数值均沿正文给定的简化口径。}

\subsection{参考答案}
\begin{enumerate}[label=\textbf{A\arabic*.},itemsep=0.45em]
\item 容量增大允许保存更多参数、状态或激活，但没有提高每秒计算或搬运速度。大 batch 可能提高复用和吞吐，也会增加每步工作量，所以单次运算或每步耗时不能由容量直接推出。
\item 元素数为 $32\times16\times16\times64=524\,288$；fp32 为 $2\,097\,152$ byte，即 2 MiB；bf16 为 1 MiB。bf16 与 fp32 都有 8 位指数，但 bf16 尾数仅 7 位，fp32 有 23 位；范围和精度是不同属性。
\item 输出 $(2,3,5)$；hidden 轴上的对应元素相乘后求和，保留 batch 与两个位置轴。沿 hidden 求均值是单张量的归约，不产生两两位置的内积矩阵，也有不同的缩放。
\item 前向约 $2\times4\times8\times8=512$ FLOPs。若上游 $G:(4,8)$，则 $GW^\top:(4,8)$，$X^\top G:(8,8)$，各约 512 FLOPs，反向合计约 1024。不需要输入梯度时，可省去 $GW^\top$，但仍要求权重梯度。
\item $N=3\times8^2=192$。参数 384、梯度 384、AdaGrad 状态 768、激活粗估 192 byte，总计 1728 byte。Adam 两份状态合计 1536 byte，总计 2496 byte；均未计输入、额外中间量与运行开销。
\item 近似 $F=2\times1024^3=2\,147\,483\,648$ FLOPs，$Q=6\times1024^2=6\,291\,456$ byte，$I\approx341.3$；若精确计点积加法则 $I\approx341.2$，均大于 295.4，在理想模型中计算受限。矩阵向量乘法约 $2n^2$ FLOPs，却仍需读约 $2n^2$ byte 权重，强度趋近 1，权重复用少。
\item 总计算量约 $6.3\times10^{24}$ FLOPs。时间为该数除以 $1024\times989.5\times10^{12}\times0.5$，再除以 86400，约 143.9 天。MFU 减半则估计时间翻倍到约 287.9 天；这一结论以其他条件固定为前提。
\item 每个大小 16 的 micro-batch 均值损失除以 4，再 backward；四次全部完成后只 step 一次，并清梯度。简化激活峰值约降至四分之一；参数、梯度、优化器状态基本不变。不等长或有 mask 时要按实际参与损失的样本／token 加权。
\item $L/k$ 项用于保存各段边界，$k$ 项用于反向时恢复并保存当前段中间量；平衡得 $k\approx\sqrt L$，激活量级为 $O(\sqrt L)$。每段重算约 $k$ 层，约 $L/k$ 段，总额外重算为 $O(L)$。它主要节省激活，不能消除参数本身的占用。
\end{enumerate}

\nextpage
\subsection{阅读范围与来源}\label{sec:sources}
\small
核心依据为 \textbf{Spring 2026 第 2 课正式可执行讲义、官网标注的录像版可执行讲义及其代码和执行记录}。本资料按问题链重新组织，不是录像逐字稿；没有把读取录像版讲义描述为观看完整视频，也不声称逐篇精读所有背景论文。

\begin{itemize}[itemsep=0.25em]
\item \href{https://cs336.stanford.edu/}{课程主页与 Schedule}：第 2 课日期为 2026 年 4 月 1 日，主题为 PyTorch（einops）与资源核算。
\item \href{https://github.com/stanford-cs336/lectures/blob/main/lecture_02.py}{正式讲义源码}；\href{https://cs336.stanford.edu/lectures/?trace=lecture_02}{正式在线讲义}；\href{https://cs336.stanford.edu/lectures/?trace=lecture_02_recording}{录像版在线讲义}：张量、精度、FLOPs、roofline、梯度、优化器及显存优化的全部主体。
\item \href{https://einops.rocks/1-einops-basics/}{einops 官方教程}：轴的组合、拆分、重排与归约。
\item \href{https://docs.pytorch.org/docs/stable/amp.html}{PyTorch AMP 文档}；\href{https://docs.pytorch.org/docs/stable/checkpoint.html}{checkpoint 文档}：补充核对 autocast 的适用范围与检查点接口约定。
\item \href{https://arxiv.org/abs/1710.03740}{Mixed Precision Training}：fp32 主权重、loss scaling 的背景；\href{https://arxiv.org/abs/2209.05433}{FP8 Formats for Deep Learning}：E4M3 与 E5M2；\href{https://developer.nvidia.com/blog/introducing-nvfp4-for-efficient-and-accurate-low-precision-inference/}{NVIDIA NVFP4 说明}：低位宽及缩放开销。只提取本课相关机制，不展开完整训练配方。
\item \href{https://www.jmlr.org/papers/volume12/duchi11a/duchi11a.pdf}{AdaGrad 原论文}：核对累计平方梯度与逐坐标调整的含义，不展开收敛证明。
\item \href{https://jax-ml.github.io/scaling-book/roofline/}{All About Rooflines}：核对时间界、强度与数据复用；\href{https://www.nvidia.com/en-us/data-center/h100/}{NVIDIA H100 官方资料}及\href{https://github.com/stanford-cs336/lectures/blob/main/facts.py}{课程 facts.py}：硬件口径。正文保留课堂采用的数值，不把它当作所有 H100 型号与精度的统一规格。
\item 讲义延伸入口：\href{https://www.adamcasson.com/posts/transformer-flops}{Transformer FLOPs}（已读取，对照 $6N$ 与 MFU 的边界）；\href{https://erees.dev/transformer-memory/}{Transformer 显存分析}（本次返回 404，未用其不可访问内容作依据）。
\end{itemize}

\subsection{版本差异与补充解释的边界}
正式讲义、部署的执行快照和录像版并非完全相同。录像版开场写 B100，代码却使用 H100 峰值；另有“2 周”标题但只乘 1 周、把参数字节数再次乘字节宽度、将梯度累积元素数称为 memory，以及检查点总重算复杂度等不一致。正文按\textbf{单位一致、公式可推导}的口径核对，不复制这些数值或文字错误。

激活字节数与 PyTorch 实际分配、AMP 与永久存储类型的区别、8 卡聚合的切分前提、MFU 的理想界、损失加权与检查点分段推导，是为复习补充的澄清。课堂训练循环只演示接口；正文用标准标量损失流程说明，不把演示数据或广播行为当作真实任务配置。

讲义还列出 Transformer 显存和 FLOPs 的延伸文章，以及模型案例、训练动态和课程行政链接。它们属于背景或后续入口；本资料不将模型架构、历史排行榜、未公开模型的估算或作业答案扩展为本课要求。所有核心章节已按正式与录像版逐项覆盖。

公开资料可能继续更新；整理日期为 2026 年 10 月 23 日。源码版本和执行快照指纹、数值及排版检查记录见同目录 README 与检查记录；原始资料通过以上官方链接访问。
\normalsize
\end{document}
